% ===== PRÉAMBULE CANONIQUE — à coller VERBATIM en tête de CHAQUE .tex =====
\documentclass[11pt,a4paper]{article}
\usepackage[utf8]{inputenc}\usepackage[T1]{fontenc}\usepackage{lmodern}
\usepackage[french]{babel}
\usepackage{amsmath,amssymb,mathtools}\usepackage{icomma}
\usepackage[version=4]{mhchem}          % \ce{...} : équations & formules
\usepackage{siunitx}\sisetup{locale=FR} % \qty{}{}, \si{}, \ang{}
\usepackage{chemfig}                    % \chemfig{...} : structures développées
\usepackage{xcolor}\usepackage{colortbl}
\usepackage{tikz}\usepackage{pgfplots}\pgfplotsset{compat=1.18}
\usepackage[most]{tcolorbox}\usepackage{enumitem}\usepackage{array,booktabs}
\usepackage[a4paper,margin=2.2cm]{geometry}
\usetikzlibrary{arrows.meta,calc,angles,quotes,positioning,patterns,backgrounds,shapes.geometric}
\definecolor{primary}{RGB}{222,49,99}\definecolor{primarydark}{RGB}{150,22,55}
\definecolor{defcol}{RGB}{0,114,178}\definecolor{methcol}{RGB}{52,92,148}
\definecolor{excol}{RGB}{0,135,90}\definecolor{attncol}{RGB}{200,40,55}
\tcbset{boxbase/.style={enhanced,breakable,boxrule=0.9pt,arc=2.5pt,
  left=3mm,right=3mm,top=2.3mm,bottom=2.3mm,fonttitle=\bfseries,coltitle=white}}
% --- boîtes TUTORIEL (noms EXACTS, ne pas renommer) ---
\newtcolorbox{cle}[1][]{boxbase,colback=defcol!6,colframe=defcol,
  title={\textsc{À retenir}\ifx\relax#1\relax\relax\else\textmd{ : \itshape #1}\fi}}
\newtcolorbox{methode}[1][]{boxbase,colback=methcol!7,colframe=methcol,sharp corners,
  title={\textsc{Méthode}\ifx\relax#1\relax\relax\else\textmd{ --- \itshape #1}\fi}}
\newtcolorbox{exemplebox}[1][]{boxbase,colback=excol!5,colframe=excol,
  title={\textsc{Exemple}\ifx\relax#1\relax\relax\else\textmd{ : \itshape #1}\fi}}
\newtcolorbox{piege}[1][]{boxbase,colback=attncol!6,colframe=attncol,
  title={\textsc{Piège}\ifx\relax#1\relax\relax\else\textmd{ : \itshape #1}\fi}}
\newtcolorbox{remarque}[1][]{boxbase,colback=black!4,colframe=black!45,
  title={\textsc{Remarque}\ifx\relax#1\relax\relax\else\textmd{ : \itshape #1}\fi}}
% --- boîtes EXERCICE / CORRIGÉ (noms EXACTS, auto-comptées) ---
\newtcolorbox[auto counter]{exo}[1][]{boxbase,colback=primary!4,colframe=primary,
  title={\textsc{Exercice}~\thetcbcounter\ifx\relax#1\relax\relax\else\textmd{ --- \itshape #1}\fi}}
\newtcolorbox[auto counter]{corrige}[1][]{boxbase,colback=excol!4,colframe=excol,
  title={\textsc{Corrigé}~\thetcbcounter\ifx\relax#1\relax\relax\else\textmd{ --- \itshape #1}\fi}}
\newcommand{\R}{\mathbb{R}}\newcommand{\N}{\mathbb{N}}\newcommand{\Z}{\mathbb{Z}}\newcommand{\ds}{\displaystyle}
% (un \usepackage{fancyhdr}+en-tête est possible ; il est ignoré côté web)
% =========================================================================
\begin{document}

\begin{center}{\large\bfseries\color{primarydark} Thème 5 — Corrigé Difficile}\end{center}

\begin{corrige}[de la solution commerciale à la solution diluée]
\begin{enumerate}[label=\alph*)]
  \item Sur $\SI{1}{\litre}$ : $m_{\text{sol}} = 1000\times 1{,}10 = \SI{1100}{\gram}$, dont
        $m(\ce{NaOH}) = 1100\times 0{,}20 = \SI{220}{\gram}$, soit $n = \dfrac{220}{40} = \SI{5.5}{\mole}$.
        Donc $C_1 = \SI{5.5}{\mole\per\litre}$.
  \item $V_1 = \dfrac{C_2 V_2}{C_1} = \dfrac{0{,}50\times 250}{5{,}5} = \SI{22.7}{\milli\litre}$.
\end{enumerate}
\end{corrige}

\begin{corrige}[dilution en fonction du volume prélevé]
\begin{enumerate}[label=\alph*)]
  \item La quantité de soluté se conserve : $n = C_0 V_1 = C_2 V_2$. Avec $V_2 = 4V_1$ :
        $C_2 = \dfrac{C_0 V_1}{4 V_1} = \dfrac{C_0}{4}$.
  \item $C_2 = \dfrac{1{,}2}{4} = \SI{0.30}{\mole\per\litre}$.
\end{enumerate}
\end{corrige}

\begin{corrige}[mélange à trois constituants]
\begin{enumerate}[label=\alph*)]
  \item $n(\ce{H2O}) = \dfrac{36}{18} = \SI{2}{\mole}$ ; $n(\text{glycérol}) = \dfrac{92}{92} = \SI{1}{\mole}$ ;
        $n(\text{éthanol}) = \dfrac{46}{46} = \SI{1}{\mole}$.
  \item $n_{\text{tot}} = \SI{4}{\mole}$ : $\chi_{\ce{H2O}} = \dfrac{2}{4} = 0{,}50$,
        $\chi_{\text{glycérol}} = \dfrac{1}{4} = 0{,}25$, $\chi_{\text{éthanol}} = \dfrac{1}{4} = 0{,}25$.
        Somme : $0{,}50 + 0{,}25 + 0{,}25 = 1$.
\end{enumerate}
\end{corrige}

\begin{corrige}[molarité \emph{ou} molalité~?]
\begin{enumerate}[label=\alph*)]
  \item $m = \dfrac{n}{m_{\text{solvant}}} = \dfrac{0{,}50}{1{,}0} = \SI{0.50}{\mole\per\kilo\gram}$.
  \item Masse de soluté : $0{,}50\times 58{,}5 = \SI{29.25}{\gram}$, donc
        $m_{\text{sol}} = 1000 + 29{,}25 = \SI{1029.25}{\gram}$.
  \item $V = \dfrac{m_{\text{sol}}}{\rho} = \dfrac{1029{,}25}{1{,}02} = \SI{1009}{\milli\litre} = \SI{1.009}{\litre}$.
  \item $C = \dfrac{n}{V} = \dfrac{0{,}50}{1{,}009} = \SI{0.496}{\mole\per\litre}$. Elle est \emph{proche}
        mais \textbf{différente} de la molalité : la molarité rapporte $n$ au \emph{volume de solution},
        la molalité à la \emph{masse de solvant}.
\end{enumerate}
\end{corrige}

\begin{corrige}[mélange de deux solutions du même soluté]
La quantité de glucose se conserve ; on somme puis on divise par le volume total ($\SI{0.500}{\litre}$).
\[ n = 0{,}200\times 0{,}30 + 0{,}300\times 0{,}80 = 0{,}060 + 0{,}240 = \SI{0.300}{\mole}. \]
\[ C = \frac{0{,}300}{0{,}500} = \SI{0.600}{\mole\per\litre}. \]
\end{corrige}

\end{document}
